\subsection{Взаимодействия с горками}
%%%%%%%%%%%%%%%%%%%%%%%%%%%%%%%%%%%%%%%%%%%%%%%%%%%%
%%%%%%%%%%%%%%%%%%%%%%%%%%%%%%%%%%%%%%%%%%%%%%%%%%%%	
\z[{\begin{wrapfigure}[4]{r}{0.3\textwidth}%
		\vspace{-10pt}%
		\begin{flushright}%
 			\includegraphics[width=0.3\textwidth]{7.5.1}%
		\end{flushright}%
	\end{wrapfigure}%
	}]%
	{По горизонтальной плоскости может скользить без трения гладкая «горка» высоты $h$ и массы $m_1$. Горка плавно переходит в плоскость. При какой наименьшей скорости горки небольшое тело массы $m_2$, неподвижно лежащее в начале на ее пути, перевалит через вершину?}
%
{$v_0 = \sqrt{2gH\left(1 + \frac{m}{M}\right)}$}
%%%%%%%%%%%%%%%%%%%%%%%%%%%%%%%%%%%%%%%%%%%%%%%%%%%% 
\z[{\begin{wrapfigure}[5]{r}{0.3\textwidth}%
		\vspace{-15pt}%
		\begin{flushright}%
 			\includegraphics[width=0.3\textwidth]{7.5.2}%
		\end{flushright}%
	\end{wrapfigure}%
	}]%
	{На гладкой горизонтальной поверхности стола покоится горка с двумя вершинами, высоты которых $H$ и $3H$. На левой вершине горки находится шайба массой $m$. Масса горки $5m$, её поверхность гладкая. От незначительного толчка вправо шайба приходит в движение. Найти скорость шайбы на правой вершине, если:
	\lpu горка закреплена на столе;\lpuend
	\lpu горка не закреплена. \lpuend
	}
%
{\puan  $v = 2\sqrt{gH};$~\puan $v = \sqrt{\frac{10}{3}gH}$}
%%%%%%%%%%%%%%%%%%%%%%%%%%%%%%%%%%%%%%%%%%%%%%%%%%%%
\z[{\begin{wrapfigure}[6]{r}{0.3\textwidth}%
		\vspace{-15pt}%
		\begin{flushright}%
 			\includegraphics[width=0.3\textwidth]{7.5.3}%
		\end{flushright}%
	\end{wrapfigure}%
	}]%
	{На вершине покоящейся на гладком горизонтальном столе горки массой $3m$ удерживают шайбу массой $m$. Шайбу отпускают, и она скользит по горке без трения и отрыва и покидает горку. Горка, не отрывавшаяся от стола, приобретает скорость $u$. Найти разность высот $H$ между вершинами горки. Верхняя часть поверхности правой вершины горки наклонена к вертикали под углом $\beta = 30\degree$.}
%
{$H = \frac{2u^2}{g}\left(3 + 4\ctg^2\beta\right) = \frac{30u^2}{g}$}
%%%%%%%%%%%%%%%%%%%%%%%%%%%%%%%%%%%%%%%%%%%%%%%%%%%%	
\z[{\begin{wrapfigure}[5]{r}{0.3\textwidth}%
		\vspace{-15pt}%
		\begin{flushright}%
 			\includegraphics[width=0.3\textwidth]{7.5.4}%
		\end{flushright}%
	\end{wrapfigure}%
	}]%
	{Брусок с выемкой в форме полуцилиндра радиусом $R$ движется со скоростью $u$ по гладкой горизонтальной поверхности стола (см. рисунок). Небольшая по сравнению с размерами бруска монета массой $m$ скользит по столу со скоростью $v$ навстречу бруску, скользит далее по гладкой поверхности выемки, не отрываясь от неё, и оказывается в точке $B$, продолжая скользить по выемке вверх. Радиус $OB$ составляет угол $\varphi$ ($\cos\varphi = \frac{2}{3}$) с вертикалью. Масса бруска намного больше массы монеты.
	\lpu Найдите скорость монеты относительно бруска в точке $B$. \lpuend
	\lpu Найдите силу давления монеты на брусок в точке $B$. \lpuend
	}
%
{\puan  $v_B = \sqrt{(v + u)^2 - \frac{10}{3}gR};$~\puan $N = m\left(\frac{(v + u)^2}{R} - 4g\right)$}
%%%%%%%%%%%%%%%%%%%%%%%%%%%%%%%%%%%%%%%%%%%%%%%%%%%% 
\z[{\begin{wrapfigure}[2]{r}{0.3\textwidth}%
		\vspace{-10pt}%
		\begin{flushright}%
 			\includegraphics[width=0.3\textwidth]{7.5.5}%
		\end{flushright}%
	\end{wrapfigure}%
	}]%
	{Небольшая шайба, скользящая по гладкой горизонтальной поверхности, наезжает на гладкую горку, покоящуюся на той же поверхности (рис.). После того как шайба соскользнула с горки, оказалось, что шайба и горка движутся по гладкой горизонтальной поверхности с одинаковыми по модулю скоростями.
	\lpu Определите, при каком соотношении масс шайбы и горки это возможно. \lpuend
	\lpu Найдите отношение максимальной потенциальной энергии, которая была у шайбы во время подъёма на горку, к начальной кинетической энергии шайбы. \lpuend
	}
%
{\puan  $\frac{1}{3};$~\puan $\frac{3}{4}$}